% Sorgente LaTeX canonico, importato e revisionato dal precedente sorgente Markdown.
\chapter*{Guida alla lettura della simbologia e delle formule}
\addcontentsline{toc}{chapter}{Guida alla lettura della simbologia e delle formule}
\markboth{Guida alla lettura della simbologia e delle formule}{Guida alla lettura della simbologia e delle formule}

\begin{obiettivocapitolo}{leggere senza ambiguità simboli, condizioni e rami}
Questa guida non è materia da ripetere a memoria. Va consultata quando una notazione non è chiara o quando una formula contiene condizioni che modificano il risultato.
\end{obiettivocapitolo}
\begin{proceduraesame}{lettura di una formula}
\begin{enumerate}
\item Individua la grandezza a sinistra: è ciò che la formula calcola o definisce.
\item Leggi quantificatori e condizioni: dominio, segni, valori esclusi e regime di limite.
\item Distingui uguaglianza, equivalenza e implicazione. Solo l'equivalenza consente di procedere in entrambi i versi.
\item Se compaiono radici, $\pm$, logaritmi o inverse trigonometriche, identifica tutti i rami ammessi.
\end{enumerate}
\end{proceduraesame}

Questa sezione introduce la notazione usata nel formulario. Le condizioni poste accanto a una formula --- dominio, segni, valori esclusi e regolarità --- sono parte integrante della formula.

\section{Come interpretare i diversi tipi di scrittura}\label{come-interpretare-i-diversi-tipi-di-scrittura}

\begin{longtable}[]{@{}
  >{\raggedright\arraybackslash}p{(\columnwidth - 4\tabcolsep) * \real{0.3333}}
  >{\raggedright\arraybackslash}p{(\columnwidth - 4\tabcolsep) * \real{0.3333}}
  >{\raggedright\arraybackslash}p{(\columnwidth - 4\tabcolsep) * \real{0.3333}}@{}}
\toprule\noalign{}
\begin{minipage}[b]{\linewidth}\raggedright
Scrittura
\end{minipage} & \begin{minipage}[b]{\linewidth}\raggedright
Si legge
\end{minipage} & \begin{minipage}[b]{\linewidth}\raggedright
Uso corretto
\end{minipage} \\
\midrule\noalign{}
\endhead
\bottomrule\noalign{}
\endlastfoot
\(A:=B\) & «\(A\) è definito come \(B\)» & introduce una definizione \\
\(A=B\) & «\(A\) è uguale a \(B\)» & uguaglianza valida nelle condizioni indicate \\
\(A\equiv B\) & «\(A\) è identicamente uguale a \(B\)» & identità valida per tutti i valori del dominio \\
\(P\Rightarrow Q\) & «\(P\) implica \(Q\)» & da \(P\) segue \(Q\); il converso non è automatico \\
\(P\Longleftrightarrow Q\) & «\(P\) se e solo se \(Q\)» & valgono sia \(P\Rightarrow Q\) sia \(Q\Rightarrow P\) \\
\(A\not\Longrightarrow B\) & «\(A\) non implica \(B\)» & segnala che il converso o la deduzione proposta è falsa \\
\(A\approx B\) & «\(A\) è approssimativamente uguale a \(B\)» & approssimazione numerica o locale, non uguaglianza esatta \\
\(\pm\) & «più o meno» & rappresenta due casi distinti, uno con \(+\) e uno con \(-\) \\
\(\mp\) & «meno o più» & usa il segno opposto a quello scelto in \(\pm\) \\
\(f\sim g\) & «\(f\) è asintoticamente equivalente a \(g\)» & \(f/g\to1\) nel regime specificato \\
\(f=o(g)\) & «\(f\) è un o-piccolo di \(g\)» & \(f/g\to0\) quando il rapporto è definito \\
\(f=O(g)\) & «\(f\) è un O-grande di \(g\)» & il modulo di \(f\) è limitato da una costante per il modulo di \(g\) \\
\end{longtable}

\section{Convenzioni generali sulle lettere}\label{convenzioni-generali-sulle-lettere}

Le condizioni locali prevalgono sempre sulle convenzioni seguenti.

\begin{longtable}[]{@{}
  >{\raggedright\arraybackslash}p{(\columnwidth - 2\tabcolsep) * \real{0.5000}}
  >{\raggedright\arraybackslash}p{(\columnwidth - 2\tabcolsep) * \real{0.5000}}@{}}
\toprule\noalign{}
\begin{minipage}[b]{\linewidth}\raggedright
Simbolo
\end{minipage} & \begin{minipage}[b]{\linewidth}\raggedright
Uso abituale nel formulario
\end{minipage} \\
\midrule\noalign{}
\endhead
\bottomrule\noalign{}
\endlastfoot
\(x,y,t,u\) & variabili reali; \(t\) e \(u\) sono spesso variabili di sostituzione \\
\(a,b,c,h,k\) & parametri, estremi o costanti; il significato è indicato localmente \\
\(n,m,k\) & indici; l'insieme ammesso è specificato localmente, di solito \(\mathbb N_0\), \(\mathbb N_+\) oppure \(\mathbb Z\) \\
\(\mathbb N\) & nel formulario coincide con \(\mathbb N_+=\{1,2,3,\ldots\}\) \\
\(\varepsilon,\delta\) & quantità positive arbitrarie usate nelle definizioni di limite e continuità \\
\(C,C_0,C_1,\ldots\) & costanti arbitrarie nelle primitive \\
\(D_f\) & dominio naturale o dominio assegnato della funzione \(f\) \\
\(\operatorname{Im}f\) & immagine di \(f\), cioè insieme dei valori effettivamente assunti \\
\(\overline{\mathbb R}\) & retta reale estesa \(\mathbb R\cup\{-\infty,+\infty\}\) \\
\(\mathbb K\) & campo scalare, di solito \(\mathbb R\) oppure \(\mathbb C\) \\
\end{longtable}

Gli angoli sono espressi in \textbf{radianti}, salvo indicazione esplicita in gradi. Nelle primitive la costante \(C\) può assumere valori differenti sulle diverse componenti connesse del dominio.

\section{Simboli logici}\label{simboli-logici}

\begin{longtable}[]{@{}
  >{\raggedright\arraybackslash}p{(\columnwidth - 4\tabcolsep) * \real{0.3333}}
  >{\raggedright\arraybackslash}p{(\columnwidth - 4\tabcolsep) * \real{0.3333}}
  >{\raggedright\arraybackslash}p{(\columnwidth - 4\tabcolsep) * \real{0.3333}}@{}}
\toprule\noalign{}
\begin{minipage}[b]{\linewidth}\raggedright
Simbolo
\end{minipage} & \begin{minipage}[b]{\linewidth}\raggedright
Si legge
\end{minipage} & \begin{minipage}[b]{\linewidth}\raggedright
Significato
\end{minipage} \\
\midrule\noalign{}
\endhead
\bottomrule\noalign{}
\endlastfoot
\(\forall\) & per ogni & la proprietà vale per tutti gli elementi considerati \\
\(\exists\) & esiste & esiste almeno un elemento che soddisfa la proprietà \\
\(\exists!\) & esiste ed è unico & esiste un solo elemento con quella proprietà \\
\(\nexists\) & non esiste & nessun elemento soddisfa la proprietà \\
\(\neg P\) & non \(P\) & negazione della proposizione \(P\) \\
\(P\land Q\) & \(P\) e \(Q\) & entrambe le proposizioni devono essere vere \\
\(P\lor Q\) & \(P\) oppure \(Q\) & è sufficiente che almeno una sia vera; l'«oppure» è inclusivo \\
\(P\Rightarrow Q\) & se \(P\), allora \(Q\) & implicazione \\
\(P\Longleftrightarrow Q\) & \(P\) se e solo se \(Q\) & equivalenza \\
\(:\) & tale che & separa una condizione dagli oggetti che la soddisfano \\
\end{longtable}

Esempio:

\[
\forall x\in A\ \exists!y\in B:\ y=f(x)
\]

si legge: «per ogni \(x\) appartenente ad \(A\) esiste un unico \(y\) appartenente a \(B\) tale che \(y=f(x)\)».

\section{Insiemi, intervalli e relazioni di appartenenza}\label{insiemi-intervalli-e-relazioni-di-appartenenza}

\begin{longtable}[]{@{}
  >{\raggedright\arraybackslash}p{(\columnwidth - 4\tabcolsep) * \real{0.3333}}
  >{\raggedright\arraybackslash}p{(\columnwidth - 4\tabcolsep) * \real{0.3333}}
  >{\raggedright\arraybackslash}p{(\columnwidth - 4\tabcolsep) * \real{0.3333}}@{}}
\toprule\noalign{}
\begin{minipage}[b]{\linewidth}\raggedright
Simbolo
\end{minipage} & \begin{minipage}[b]{\linewidth}\raggedright
Si legge
\end{minipage} & \begin{minipage}[b]{\linewidth}\raggedright
Significato
\end{minipage} \\
\midrule\noalign{}
\endhead
\bottomrule\noalign{}
\endlastfoot
\(x\in A\) & \(x\) appartiene ad \(A\) & \(x\) è un elemento dell'insieme \(A\) \\
\(x\notin A\) & \(x\) non appartiene ad \(A\) & \(x\) non è elemento di \(A\) \\
\(A\subseteq B\) & \(A\) è contenuto in \(B\) & ogni elemento di \(A\) appartiene a \(B\); è ammesso \(A=B\) \\
\(A\subset B\) & \(A\) è contenuto propriamente in \(B\) & \(A\subseteq B\) e \(A\ne B\) \\
\(\varnothing\) & insieme vuoto & insieme privo di elementi \\
\(A\cup B\) & unione di \(A\) e \(B\) & elementi che appartengono ad almeno uno dei due insiemi \\
\(A\cap B\) & intersezione di \(A\) e \(B\) & elementi comuni ad \(A\) e \(B\) \\
\(A\setminus B\) & \(A\) meno \(B\) & elementi di \(A\) che non appartengono a \(B\) \\
\(A^c\) & complementare di \(A\) & elementi dell'insieme universo che non appartengono ad \(A\) \\
\(\mathcal P(A)\) & insieme delle parti di \(A\) & insieme di tutti i sottoinsiemi di \(A\) \\
\(\operatorname{card}A\) & cardinalità di \(A\) & numero di elementi, o grandezza insiemistica, di \(A\) \\
\(\{x:P(x)\}\) & insieme degli \(x\) tali che \(P(x)\) & notazione per proprietà caratteristica \\
\end{longtable}

Intervalli reali:

\begin{longtable}[]{@{}
  >{\raggedright\arraybackslash}p{(\columnwidth - 4\tabcolsep) * \real{0.3333}}
  >{\raggedright\arraybackslash}p{(\columnwidth - 4\tabcolsep) * \real{0.3333}}
  >{\raggedright\arraybackslash}p{(\columnwidth - 4\tabcolsep) * \real{0.3333}}@{}}
\toprule\noalign{}
\begin{minipage}[b]{\linewidth}\raggedright
Scrittura
\end{minipage} & \begin{minipage}[b]{\linewidth}\raggedright
Si legge
\end{minipage} & \begin{minipage}[b]{\linewidth}\raggedright
Estremi inclusi
\end{minipage} \\
\midrule\noalign{}
\endhead
\bottomrule\noalign{}
\endlastfoot
\((a,b)\) & intervallo aperto da \(a\) a \(b\) & nessuno \\
\([a,b]\) & intervallo chiuso da \(a\) a \(b\) & entrambi \\
\((a,b]\) & aperto a sinistra e chiuso a destra & soltanto \(b\) \\
\([a,b)\) & chiuso a sinistra e aperto a destra & soltanto \(a\) \\
\((a,+\infty)\) & reali maggiori di \(a\) & \(+\infty\) non è mai un estremo incluso \\
\end{longtable}

Gli oggetti topologici più frequenti sono:

\begin{longtable}[]{@{}
  >{\raggedright\arraybackslash}p{(\columnwidth - 4\tabcolsep) * \real{0.3333}}
  >{\raggedright\arraybackslash}p{(\columnwidth - 4\tabcolsep) * \real{0.3333}}
  >{\raggedright\arraybackslash}p{(\columnwidth - 4\tabcolsep) * \real{0.3333}}@{}}
\toprule\noalign{}
\begin{minipage}[b]{\linewidth}\raggedright
Simbolo
\end{minipage} & \begin{minipage}[b]{\linewidth}\raggedright
Si legge
\end{minipage} & \begin{minipage}[b]{\linewidth}\raggedright
Significato
\end{minipage} \\
\midrule\noalign{}
\endhead
\bottomrule\noalign{}
\endlastfoot
\(I(x_0,r)\) & intorno di centro \(x_0\) e raggio \(r\) & \((x_0-r,x_0+r)\) \\
\(I^*(x_0,r)\) & intorno puntato & intorno dal quale è escluso il centro \(x_0\) \\
\(A^\circ\) o \(\operatorname{int}A\) & interno di \(A\) & insieme dei punti interni \\
\(A'\) & derivato di \(A\) & insieme dei punti di accumulazione di \(A\) \\
\(\overline A\) & chiusura di \(A\) & \(A\) unito ai suoi punti di accumulazione \\
\(\partial A\) & frontiera di \(A\) & punti di confine tra \(A\) e il suo complementare \\
\end{longtable}

\section{Insiemi numerici e operatori discreti}\label{insiemi-numerici-e-operatori-discreti}

\begin{longtable}[]{@{}
  >{\raggedright\arraybackslash}p{(\columnwidth - 4\tabcolsep) * \real{0.3333}}
  >{\raggedright\arraybackslash}p{(\columnwidth - 4\tabcolsep) * \real{0.3333}}
  >{\raggedright\arraybackslash}p{(\columnwidth - 4\tabcolsep) * \real{0.3333}}@{}}
\toprule\noalign{}
\begin{minipage}[b]{\linewidth}\raggedright
Simbolo
\end{minipage} & \begin{minipage}[b]{\linewidth}\raggedright
Si legge
\end{minipage} & \begin{minipage}[b]{\linewidth}\raggedright
Insieme o operatore
\end{minipage} \\
\midrule\noalign{}
\endhead
\bottomrule\noalign{}
\endlastfoot
\(\mathbb N_0\) & naturali incluso lo zero & \(\{0,1,2,\ldots\}\) \\
\(\mathbb N\) o \(\mathbb N_+\) & naturali positivi & \(\{1,2,3,\ldots\}\) \\
\(\mathbb Z\) & interi relativi & \(\ldots,-2,-1,0,1,2,\ldots\) \\
\(\mathbb Z_{<0}\) & interi negativi & \(\{-1,-2,-3,\ldots\}\) \\
\(\mathbb Q\) & razionali & numeri rappresentabili come rapporto di interi \\
\(\mathbb R\) & reali & retta reale \\
\(\mathbb C\) & complessi & numeri della forma \(a+ib\) \\
\(\sum_{k=m}^{n}a_k\) & sommatoria da \(k=m\) a \(k=n\) & \(a_m+a_{m+1}+\cdots+a_n\) \\
\(\prod_{k=m}^{n}a_k\) & produttoria da \(k=m\) a \(k=n\) & \(a_ma_{m+1}\cdots a_n\) \\
\(n!\) & fattoriale di \(n\) & prodotto degli interi da \(1\) a \(n\) \\
\(\binom nk\) & coefficiente binomiale \(n\) su \(k\) & numero di scelte di \(k\) elementi fra \(n\) \\
\(a\mid b\) & \(a\) divide \(b\) & esiste un intero \(k\) tale che \(b=ak\) \\
\(n\bmod m\) & resto di \(n\) modulo \(m\) & resto della divisione euclidea, con \(m>0\) \\
\(a\equiv b\pmod m\) & \(a\) congruo a \(b\) modulo \(m\) & \(m\) divide \(a-b\), con \(m\in\mathbb N_+\) \\
\(\lfloor x\rfloor\) & parte intera inferiore di \(x\) & massimo intero non maggiore di \(x\) \\
\(\operatorname{sgn}x\) & segno di \(x\) & vale \(-1\), \(0\) oppure \(1\) secondo il segno di \(x\) \\
\(\gcd(a,b)\) & massimo comune divisore & massimo divisore positivo comune ad \(a\) e \(b\) \\
\(\operatorname{lcm}(a,b)\) & minimo comune multiplo & minimo multiplo positivo comune ad \(a\) e \(b\) \\
\(\deg P\) & grado di \(P\) & massimo esponente con coefficiente non nullo \\
\end{longtable}

\section{Matrici e sistemi lineari}\label{matrici-e-sistemi-lineari}

\begin{longtable}[]{@{}
  >{\raggedright\arraybackslash}p{(\columnwidth - 4\tabcolsep) * \real{0.3333}}
  >{\raggedright\arraybackslash}p{(\columnwidth - 4\tabcolsep) * \real{0.3333}}
  >{\raggedright\arraybackslash}p{(\columnwidth - 4\tabcolsep) * \real{0.3333}}@{}}
\toprule\noalign{}
\begin{minipage}[b]{\linewidth}\raggedright
Scrittura
\end{minipage} & \begin{minipage}[b]{\linewidth}\raggedright
Si legge
\end{minipage} & \begin{minipage}[b]{\linewidth}\raggedright
Significato
\end{minipage} \\
\midrule\noalign{}
\endhead
\bottomrule\noalign{}
\endlastfoot
\(A=(a_{ij})\) & matrice \(A\) di elementi \(a_{ij}\) & tabella rettangolare di coefficienti \\
\([A\mid c]\) & matrice completa o aumentata & matrice dei coefficienti con la colonna dei termini noti \\
\(\operatorname{rank}A\) & rango di \(A\) & massimo numero di righe o colonne linearmente indipendenti \\
\(\det A\) & determinante di \(A\) & scalare associato a una matrice quadrata \\
\(A^{-1}\) & matrice inversa & soddisfa \(AA^{-1}=A^{-1}A=I\), quando esiste \\
\(I\) & matrice identità & elemento neutro del prodotto tra matrici \\
\end{longtable}

\section{Funzioni, composizione e inversa}\label{funzioni-composizione-e-inversa}

\begin{longtable}[]{@{}
  >{\raggedright\arraybackslash}p{(\columnwidth - 4\tabcolsep) * \real{0.3333}}
  >{\raggedright\arraybackslash}p{(\columnwidth - 4\tabcolsep) * \real{0.3333}}
  >{\raggedright\arraybackslash}p{(\columnwidth - 4\tabcolsep) * \real{0.3333}}@{}}
\toprule\noalign{}
\begin{minipage}[b]{\linewidth}\raggedright
Scrittura
\end{minipage} & \begin{minipage}[b]{\linewidth}\raggedright
Si legge
\end{minipage} & \begin{minipage}[b]{\linewidth}\raggedright
Significato
\end{minipage} \\
\midrule\noalign{}
\endhead
\bottomrule\noalign{}
\endlastfoot
\(f:A\to B\) & \(f\) da \(A\) in \(B\) & \(A\) è il dominio assegnato e \(B\) il codominio \\
\(x\mapsto f(x)\) & \(x\) viene mandato in \(f(x)\) & legge di corrispondenza \\
\(f(A)\) & immagine di \(A\) tramite \(f\) & insieme dei valori \(f(x)\) con \(x\in A\) \\
\(f^{-1}(Y)\) & preimmagine di \(Y\) & insieme degli \(x\) tali che \(f(x)\in Y\) \\
\(f^{-1}\) & funzione inversa di \(f\) & esiste come funzione soltanto se \(f\) è biiettiva \\
\((g\circ f)(x)\) & \(g\) composto con \(f\) & si applica prima \(f\), poi \(g\): \(g(f(x))\) \\
\(\operatorname{id}_A\) & identità su \(A\) & funzione \(x\mapsto x\) \\
\(G_f\) & grafico di \(f\) & insieme delle coppie \((x,f(x))\) \\
\(C(I)\) & funzioni continue su \(I\) & spazio delle funzioni continue \\
\(C^k(I)\) & funzioni di classe \(C^k\) & derivate continue fino all'ordine \(k\) \\
\(\mathcal R([a,b])\) & funzioni Riemann-integrabili & funzioni integrabili secondo Riemann su \([a,b]\) \\
\(\mathcal R_{\mathrm{loc}}(I)\) & localmente Riemann-integrabili & integrabili su ogni sottointervallo compatto che evita le singolarità \\
\end{longtable}

\section{Successioni, limiti e confronto asintotico}\label{successioni-limiti-e-confronto-asintotico}

\begin{longtable}[]{@{}
  >{\raggedright\arraybackslash}p{(\columnwidth - 4\tabcolsep) * \real{0.3333}}
  >{\raggedright\arraybackslash}p{(\columnwidth - 4\tabcolsep) * \real{0.3333}}
  >{\raggedright\arraybackslash}p{(\columnwidth - 4\tabcolsep) * \real{0.3333}}@{}}
\toprule\noalign{}
\begin{minipage}[b]{\linewidth}\raggedright
Scrittura
\end{minipage} & \begin{minipage}[b]{\linewidth}\raggedright
Si legge
\end{minipage} & \begin{minipage}[b]{\linewidth}\raggedright
Significato
\end{minipage} \\
\midrule\noalign{}
\endhead
\bottomrule\noalign{}
\endlastfoot
\((a_n)\) & successione \(a_n\) & funzione definita sugli indici naturali \\
\(a_n\to L\) & \(a_n\) tende a \(L\) & i termini si avvicinano arbitrariamente a \(L\) \\
\(n\to\infty\) & \(n\) tende all'infinito & si considerano indici sempre più grandi \\
\(x\to x_0^-\) & \(x\) tende a \(x_0\) da sinistra & \(x<x_0\) \\
\(x\to x_0^+\) & \(x\) tende a \(x_0\) da destra & \(x>x_0\) \\
\(x\to\pm\infty\) & \(x\) tende a più o meno infinito & si studia il comportamento per modulo grande \\
\(\limsup a_n\) & limite superiore & massimo valore di aderenza per successioni limitate \\
\(\liminf a_n\) & limite inferiore & minimo valore di aderenza per successioni limitate \\
\(a_n\nearrow\) & successione non decrescente & i termini non diminuiscono al crescere di \(n\) \\
\(a_n\searrow\) & successione non crescente & i termini non aumentano al crescere di \(n\) \\
\(f\ll g\) & \(f\) è trascurabile rispetto a \(g\) & abbreviazione di \(f/g\to0\) \\
\end{longtable}

«Definitivamente» significa «da un certo indice in poi» per le successioni e «in un opportuno intorno puntato del regime considerato» per i limiti di funzione.

La definizione

\[
\forall\varepsilon>0\ \exists\delta>0:\
0<|x-x_0|<\delta
\Rightarrow
|f(x)-L|<\varepsilon
\]

si legge: «per ogni tolleranza positiva \(\varepsilon\), esiste una distanza positiva \(\delta\) tale che, quando \(x\) è diverso da \(x_0\) ma dista da \(x_0\) meno di \(\delta\), il valore \(f(x)\) dista da \(L\) meno di \(\varepsilon\)».

\section{Derivate, differenziali e integrali}\label{derivate-differenziali-e-integrali}

\begin{longtable}[]{@{}
  >{\raggedright\arraybackslash}p{(\columnwidth - 4\tabcolsep) * \real{0.3333}}
  >{\raggedright\arraybackslash}p{(\columnwidth - 4\tabcolsep) * \real{0.3333}}
  >{\raggedright\arraybackslash}p{(\columnwidth - 4\tabcolsep) * \real{0.3333}}@{}}
\toprule\noalign{}
\begin{minipage}[b]{\linewidth}\raggedright
Scrittura
\end{minipage} & \begin{minipage}[b]{\linewidth}\raggedright
Si legge
\end{minipage} & \begin{minipage}[b]{\linewidth}\raggedright
Significato
\end{minipage} \\
\midrule\noalign{}
\endhead
\bottomrule\noalign{}
\endlastfoot
\(f'(x)\) & derivata prima di \(f\) in \(x\) & tasso di variazione istantaneo e pendenza della tangente \\
\(f^{(n)}(x)\) & derivata \(n\)-esima & derivata applicata \(n\) volte \\
\(f'_-(x_0)\), \(f'_+(x_0)\) & derivate sinistra e destra & limiti laterali del rapporto incrementale \\
\(df_{x_0}(h)\) & differenziale di \(f\) in \(x_0\) applicato a \(h\) & parte lineare dell'incremento di \(f\) \\
\(F_x\), \(F_y\) & derivate parziali di \(F\) & derivazione rispetto a \(x\) o a \(y\) mantenendo fissa l'altra variabile \\
\(\int f(x)\,dx\) & integrale indefinito di \(f\) & famiglia delle primitive \(F+C\) \\
\(\int_a^b f(x)\,dx\) & integrale di \(f\) da \(a\) a \(b\) & accumulazione o area orientata \\
\(dx\) & differenziale della variabile \(x\) & indica la variabile rispetto alla quale si integra \\
\([F(x)]_a^b\) & \(F\) calcolata tra \(a\) e \(b\) & \(F(b)-F(a)\) \\
\(\operatorname{PV}\int\) & valore principale di Cauchy & limite simmetrico; non coincide necessariamente con la convergenza impropria \\
\(\lVert f\rVert_{\infty,E}\) & norma uniforme o del supremo & \(\sup_{x\in E}\lvert f(x)\rvert\), per \(E\ne\varnothing\) e \(f:E\to\mathbb R\) limitata \\
\end{longtable}

La derivata

\[
f'(x_0)=\lim_{h\to0}\frac{f(x_0+h)-f(x_0)}h
\]

si legge: «la derivata di \(f\) in \(x_0\) è il limite, per l'incremento \(h\) che tende a zero, del rapporto tra l'incremento della funzione e l'incremento della variabile».

Nelle ipotesi della formula di Newton--Leibniz, l'integrale

\[
\int_a^b f(x)\,dx=F(b)-F(a),
\qquad F'=f,
\]

si legge: «l'integrale di \(f\) da \(a\) a \(b\) è la differenza fra i valori, in \(b\) e in \(a\), di una qualsiasi primitiva \(F\) di \(f\)».

\section{Serie numeriche e serie di funzioni}\label{serie-numeriche-e-serie-di-funzioni}

\begin{longtable}[]{@{}
  >{\raggedright\arraybackslash}p{(\columnwidth - 4\tabcolsep) * \real{0.3333}}
  >{\raggedright\arraybackslash}p{(\columnwidth - 4\tabcolsep) * \real{0.3333}}
  >{\raggedright\arraybackslash}p{(\columnwidth - 4\tabcolsep) * \real{0.3333}}@{}}
\toprule\noalign{}
\begin{minipage}[b]{\linewidth}\raggedright
Scrittura
\end{minipage} & \begin{minipage}[b]{\linewidth}\raggedright
Si legge
\end{minipage} & \begin{minipage}[b]{\linewidth}\raggedright
Significato
\end{minipage} \\
\midrule\noalign{}
\endhead
\bottomrule\noalign{}
\endlastfoot
\(\sum_{n=n_0}^{\infty}a_n\) & serie degli \(a_n\) & limite delle somme parziali, non una somma eseguita ``tutta insieme'' \\
\(s_N=\sum_{n=n_0}^{N}a_n\) & somma parziale \(N\)-esima & somma finita dei termini fino a \(N\) \\
\(R_N=\sum_{n=N+1}^{\infty}a_n\) & resto dopo il termine \(N\) & quando la serie converge, errore commesso sostituendola con \(s_N\) \\
\(f_n\xrightarrow[]{p}f\) & convergenza puntuale & la soglia può dipendere dal punto \(x\) \\
\(f_n\xrightarrow[]{u}f\) & convergenza uniforme & una sola soglia funziona per tutti i punti del dominio \\
\end{longtable}

La scrittura

\[
\sum_{n=n_0}^{\infty}a_n=A
\Longleftrightarrow
\lim_{N\to\infty}\sum_{n=n_0}^{N}a_n=A
\]

significa che la successione delle somme parziali converge al numero \(A\).

\section{Numeri complessi}\label{numeri-complessi}

\begin{longtable}[]{@{}
  >{\raggedright\arraybackslash}p{(\columnwidth - 4\tabcolsep) * \real{0.3333}}
  >{\raggedright\arraybackslash}p{(\columnwidth - 4\tabcolsep) * \real{0.3333}}
  >{\raggedright\arraybackslash}p{(\columnwidth - 4\tabcolsep) * \real{0.3333}}@{}}
\toprule\noalign{}
\begin{minipage}[b]{\linewidth}\raggedright
Simbolo
\end{minipage} & \begin{minipage}[b]{\linewidth}\raggedright
Si legge
\end{minipage} & \begin{minipage}[b]{\linewidth}\raggedright
Significato
\end{minipage} \\
\midrule\noalign{}
\endhead
\bottomrule\noalign{}
\endlastfoot
\(i\) & unità immaginaria & \(i^2=-1\) \\
\(\operatorname{Re}z\) & parte reale di \(z\) & coefficiente reale di \(z=a+ib\) \\
\(\operatorname{Im}z\) & parte immaginaria di \(z\) & coefficiente \(b\) di \(i\) \\
\(\overline z\) & coniugato di \(z\) & se \(z=a+ib\), allora \(\overline z=a-ib\) \\
\(\lvert z\rvert\) & modulo di \(z\) & distanza di \(z\) dall'origine nel piano complesso \\
\(\arg z\) & argomenti di \(z\ne0\) & insieme multivalore degli angoli modulo \(2\pi\) \\
\(\operatorname{Arg}z\) & argomento principale di \(z\ne0\) & rappresentante scelto in un intervallo prefissato \\
\end{longtable}

\section{Formule dirette, formule inverse e rami}\label{formule-dirette-formule-inverse-e-rami}

Una \textbf{formula diretta} è risolta rispetto alla grandezza che si vuole calcolare. Una \textbf{formula inversa} è la stessa relazione risolta rispetto a un'altra grandezza. L'isolamento algebrico è lecito soltanto se non si divide per zero e se si rispettano dominio, segno e numero dei rami.

Esempio multivalore:

\[
y=x^2
\Longleftrightarrow
x=\pm\sqrt y,
\qquad y\ge0.
\]

La relazione non definisce una funzione inversa globale di \(x\mapsto x^2\). Restringendo il dominio a \([0,+\infty)\) si ottiene \(x=\sqrt y\); restringendolo a \((-\infty,0]\) si ottiene \(x=-\sqrt y\).

Analogamente,

\[
\sin^2x+\cos^2x=1
\]

è un'identità, non una formula con una sola inversa. Da essa si ricava

\[
\sin x=\pm\sqrt{1-\cos^2x},
\]

ma il segno dipende dal quadrante. Quando compaiono \(\pm\), radici, logaritmi, funzioni goniometriche inverse o elevamenti a potenza, il controllo dei rami è obbligatorio.
